\chapter{NumPy (Iris — session 8)}

\section{Introduction}
In session 6 and session 7, we focused on \textbf{classes} and \textbf{inheritance} to build reusable classifiers for the Iris classification task.
This session, we introduce \textbf{NumPy}, the most important Python library for fast numerical computing. NumPy is designed to work efficiently with large arrays of numbers, which is exactly what we need for machine learning.

Recall from session 7 that our classifier interface was \texttt{fit()}, \texttt{predict\_one()}, \texttt{predict\_all()}, and \texttt{evaluate()}.
In this session, we keep that interface but implement NumPy-based data operations underneath it.

\begin{tcolorbox}[title=You will work in]
Use the following files for this session:
\begin{itemize}
    \item \texttt{session8/session8\_iris\_template.py} (main student file)
    \item \texttt{session8/iris\_session8\_utils.py} (data loading and utilities)
    \item \texttt{session8/iris\_session8\_models.py} (NumPy-based classifier model)
\end{itemize}
\textbf{How to run the file:} Run from the project root: \\
\texttt{python3 session8/session8\_iris\_template.py}
\end{tcolorbox}

\subsection{What You Will Learn}
By the end of this tutorial, you will be able to:
\begin{itemize}\setlength{\itemsep}{0pt}
\item create NumPy arrays and understand their advantages over Python lists for numerical tasks,
\item use indexing and slicing on arrays,
\item compute means with \texttt{np.mean(..., axis=\dots)},
\item use boolean masks to select rows in a dataset,
\item compute Iris class centroids and use them for classification intuition,
\item import custom functions/classes from one Python file to another.
\end{itemize}

\subsection{Tools and Setup}
\noindent
You need:
\begin{itemize}\setlength{\itemsep}{0pt}
\item Python 3,
\item NumPy installed (\texttt{pip install numpy}),
\item an \texttt{iris.csv} file.
\end{itemize}

\noindent
\textbf{Assumption about \texttt{iris.csv}:} it contains 5 columns:
\begin{center}
    \texttt{sepal\_length,sepal\_width,petal\_length,petal\_width,species}
\end{center}
If your file has a header row, we will handle it by skipping the first row.

\section{Understanding X and y}
In machine learning, data is usually separated into features and labels. 
\begin{itemize}
    \item \textbf{X:} A 2D NumPy array of shape \texttt{(n\_samples, 4)} containing the numeric features (sepal length, sepal width, petal length, petal width).
    \item \textbf{y:} A 1D NumPy array of shape \texttt{(n\_samples,)} containing the string species labels.
\end{itemize}

\newpage
\subsection{Activity 1: Load Iris into NumPy arrays (\texttt{X} and \texttt{y})}

\noindent
Create a new file called \texttt{session8/iris\_session8\_utils.py}. Start with:

\begin{minted}[fontsize=\small, linenos, breaklines=true, frame=lines]{python}
import numpy as np

def load_iris_csv(filepath="iris.csv", skip_header=1):
    """
    Loads iris.csv into:
      X: (n_samples, 4) float array
      y: (n_samples,) string array
    """
    # Use genfromtxt for mixed types (numbers + strings)
    data = np.genfromtxt(
        filepath,
        delimiter=",",
        dtype=None,          # allow mixed dtype
        encoding="utf-8",
        skip_header=skip_header
    )

    # data will be a structured array with 5 columns
    X = np.column_stack([data[i] for i in range(4)]).astype(float)
    y = np.array(data[4], dtype=str)

    return X, y
\end{minted}

Now in \texttt{session8/session8\_iris\_template.py}, write:
\begin{minted}[fontsize=\small, linenos, breaklines=true, frame=lines]{python}
import numpy as np
from session8.iris_session8_utils import load_iris_csv

X, y = load_iris_csv("iris.csv")
print("X shape:", X.shape)
print("X dtype:", X.dtype)
print("y shape:", y.shape)
print("y dtype:", y.dtype)
\end{minted}

\begin{tcolorbox}[title=Checkpoint]
Run \texttt{session8\_iris\_template.py}. Confirm you see \texttt{X shape: (150, 4)} and \texttt{y shape: (150,)}. Check that \texttt{X dtype} is \texttt{float64} and \texttt{y} is \texttt{<U10} (string).
\end{tcolorbox}

\newpage
\subsection{Activity 2: Quick NumPy warm-up (indexing and slicing)}

\noindent
Add the following to \texttt{session8\_iris\_template.py} below your prints:

\begin{minted}[fontsize=\small, linenos, breaklines=true, frame=lines]{python}
# First sample (row 0)
print("First row of X:", X[0])
print("First label y[0]:", y[0])

# First 5 rows
print("First 5 rows of X:\n", X[:5])

# Only petal features (columns 2 and 3)
petal = X[:, 2:4]
print("Petal matrix shape:", petal.shape)

# Sepal length column only (column 0)
sepal_length = X[:, 0]
print("Sepal length shape:", sepal_length.shape)
\end{minted}

\begin{tcolorbox}[title=Do: Your Task]
\begin{enumerate}
    \item What will \texttt{X[-1]} return? Add a print statement to output the last row using negative indexing.
    \item Print the first 10 labels from \texttt{y}.
\end{enumerate}
\end{tcolorbox}

\begin{tcolorbox}[title=Checkpoint]
Verify your last row matches the bottom of your CSV and your first 10 labels are all \texttt{setosa}.
\end{tcolorbox}

\newpage
\subsection{Activity 3: Compute centroids (mean feature vector) per species}

A \textbf{boolean mask} is an array of True/False values that you can use to select specific rows from another array. For example, \texttt{y == "setosa"} returns an array of booleans. We can then do \texttt{X[y == "setosa"]} to grab only the rows in \texttt{X} where the species is setosa.

\noindent
A \textbf{centroid} is the mean of the features for a class. For each species, we want:
\[
    \mu_{\text{class}} = \texttt{mean}(X_{\text{class}}, \text{axis}=0)
\]
This gives a 4D vector:
\texttt{[mean\_sepal\_length, mean\_sepal\_width, mean\_petal\_length, mean\_petal\_width]}.

\subsubsection{Step 1: Find unique species}
\begin{minted}[fontsize=\small, linenos, breaklines=true, frame=lines]{python}
classes = np.unique(y)
print("Classes:", classes)
\end{minted}

\subsubsection{Step 2: Compute centroid with a boolean mask}
\begin{minted}[fontsize=\small, linenos, breaklines=true, frame=lines]{python}
centroids = {}

for c in classes:
    mask = (y == c)          # boolean array, True where class matches
    X_c = X[mask]            # select rows belonging to class c
    mu = np.mean(X_c, axis=0)
    centroids[c] = mu

print("Centroids:")
for c in classes:
    print(c, "->", np.round(centroids[c], 2))
\end{minted}

\begin{tcolorbox}[title=Checkpoint]
Verify that the \texttt{centroids} dictionary contains three entries (one per species) and each entry is an array of 4 rounded numbers.
\end{tcolorbox}

\newpage
\subsection{Activity 4: Why \texttt{axis=0} matters}

\noindent
In NumPy, \texttt{axis=0} means: "collapse rows (samples) and compute the mean per column (feature)". \texttt{axis=1} means "collapse columns and compute the mean per row (sample)".

\begin{minted}[fontsize=\small, linenos, breaklines=true, frame=lines]{python}
example = X[:3]  # first 3 samples
print("example shape:", example.shape)
print("mean axis=0 (per feature):", np.mean(example, axis=0))
print("mean axis=1 (per sample):", np.mean(example, axis=1))
\end{minted}

\begin{tcolorbox}[title=Do: Your Task]
\begin{enumerate}
    \item In one sentence (as a comment in your code), explain the difference between \texttt{axis=0} and \texttt{axis=1}.
\end{enumerate}
\end{tcolorbox}

\newpage

\section{Troubleshooting Imports and CSV Loading}
\begin{tcolorbox}[title=Common Mistakes \& Troubleshooting]
\begin{itemize}
    \item \textbf{ModuleNotFoundError:} If Python cannot find \texttt{session8.iris\_session8\_utils}, make sure you are running the script from the project root directory, not from inside the \texttt{session8} folder.
    \item \textbf{FileNotFoundError or ValueError (CSV):} If \texttt{load\_iris\_csv} fails, check that \texttt{iris.csv} exists. If you get a ValueError about parsing strings, ensure \texttt{skip\_header=1} matches your CSV file.
\end{itemize}
\end{tcolorbox}

\subsection{Activity 5: Code organization across files (imports)}

\noindent
As projects grow, we avoid putting everything in one file. Practice splitting code into three files: \texttt{iris\_session8\_utils.py}, \texttt{iris\_session8\_models.py}, and \texttt{session8\_iris\_template.py}.

Recall from session 7 that the classifier interface should stay stable. In \texttt{iris\_session8\_models.py}, we define a \texttt{NumpyCentroidClassifier} class that keeps the same method names: \texttt{fit()}, \texttt{predict\_one()}, \texttt{predict\_all()}, and \texttt{evaluate()}.

\begin{tcolorbox}[title=Do: Your Task]
\begin{enumerate}
    \item Run \texttt{session8/session8\_iris\_template.py} and confirm the imports work and the centroid classifier correctly evaluates the dataset.
    \item Intentionally rename \texttt{session8/iris\_session8\_utils.py} to something else and observe the import error.
    \item Rename it back correctly.
\end{enumerate}
\end{tcolorbox}

\noindent
This file/module structure is important because Session 9 will import these exact same utilities and model interface for plotting-based analysis.

\newpage
\subsection{Optional Activity: Use an LLM to explain NumPy-heavy code}

\noindent
Professional developers often paste unfamiliar code into an LLM and ask for a guided explanation.
This is a \textbf{learning tool}, not a replacement for understanding.

\subsubsection{Prompt 1: Explain NumPy usage}
Copy your centroid computation code and ask:
\begin{quote}
    Explain every place NumPy is used in this script. For each NumPy function, tell me: (1) what it does, (2) what shape the input is, (3) what shape the output is.
\end{quote}

\subsubsection{Prompt 2: Explain boolean masks + axis mean}
Copy the loop and ask:
\begin{quote}
    Explain line-by-line how boolean masking works in NumPy here, and why axis=0 produces a centroid vector.
\end{quote}

\begin{tcolorbox}[title=Do: Reflection]
Write 3--5 lines (as comments in your Python file):
\begin{itemize}\setlength{\itemsep}{0pt}
\item One thing the LLM explained clearly.
\item One thing you still find confusing.
\end{itemize}
\end{tcolorbox}

\section{Summary and Reflection}

\textbf{Summary:}
\begin{itemize}
    \item \textbf{NumPy Arrays:} Provided a more efficient, faster way to store numerical data than Python lists. We explored shapes and dtypes.
    \item \textbf{Boolean Masking:} Allowed us to filter data without writing explicit \texttt{for} loops.
    \item \textbf{Axes:} \texttt{axis=0} computes operations across rows (giving per-column results), while \texttt{axis=1} computes across columns.
    \item \textbf{Modularization:} Code was successfully split across utilities, models, and execution scripts.
\end{itemize}

\textbf{Reflection Questions:}
\begin{enumerate}
    \item Why does NumPy make numerical operations faster than Python lists?
    \item What does \texttt{axis=0} mean in the context of computing centroids for the Iris dataset?
\end{enumerate}

\subsection{Submission Requirement (Graded)}

\noindent
Submit:
\begin{enumerate}
    \item \texttt{session8/session8\_iris\_template.py}
    \item \texttt{session8/iris\_session8\_utils.py}
    \item \texttt{session8/iris\_session8\_models.py}
    \item A screenshot/PDF of your terminal output showing \texttt{X shape} and \texttt{X dtype}, and printed centroids.
\end{enumerate}
